\subsection{Fundamentos te\-ó\-ri\-cos}
\label{sec:theoretical_background}

La Inteligencia Artificial (IA) constituye un área de la informática dedicada a la construcción de sistemas capaces de percibir entornos específicos, procesarlos mediante razonamiento abstracto y ejecutar acciones determinadas para lograr metas específicas de forma inteligente \cite{ref33}. El aprendizaje automático (ML) representa una subdisciplina de la IA que habilita a los sistemas computacionales para adquirir conocimiento a partir de datos y mejorar su capacidad de desempeño sin necesidad de programación explícita para cada tarea individual \cite{ref33}. Aunque estos términos frecuentemente se utilizan de forma equivalente, la IA engloba diversos enfoques para construir sistemas dotados de inteligencia, mientras que el aprendizaje automático se concentra específicamente en la adquisición de conocimiento mediante experiencia y datos \cite{ref33}. Esta relación posiciona al aprendizaje automático como uno de los métodos fundamentales en la construcción de aplicaciones contemporáneas de inteligencia artificial \cite{ref33}.

El aprendizaje profundo o deep learning constituye una rama especializada del aprendizaje automático que utiliza arquitecturas de redes neuronales multicapa para aprender representaciones jerárquicas y complejas de los datos \cite{ref34}. A diferencia de metodologías convencionales, esta aproximación extrae de forma completamente autónoma los atributos más significativos sin requerir ingeniería manual de características \cite{ref34}. Debido a esta capacidad distintiva, el aprendizaje profundo encuentra aplicación en tareas tan variadas como reconocimiento visual de imágenes, síntesis de voz, procesamiento de lenguaje natural y análisis computacional de escenas visuales \cite{ref34}. Su rendimiento sobresaliente y aptitud para procesar volúmenes masivos de información han posicionado al aprendizaje profundo como una de las tecnologías más influyentes en la inteligencia artificial contemporánea \cite{ref34}.

El reconocimiento y categorización de imágenes en contexto médico comprende la identificación de estructuras anatómicas, tipos de tejido o patologías presentes en representaciones visuales obtenidas mediante modalidades diagnósticas tales como resonancia magnética, tomografía computarizada y radiografía convencional \cite{ref35}. Para ello, los algoritmos de aprendizaje profundo analizan de forma automática las imágenes y extraen atributos significativos que permiten asignar una categoría diagnóstica apropiada \cite{ref35}. Este proceso facilita la visualización de enfermedades, la discriminación entre tejidos fisiológicos y afecciones patológicas, y el apoyo a procesos decisorios en contextos clínicos \cite{ref35}. Gracias a su exactitud y potencia para procesar volúmenes considerables de información visual, esta metodología constituye una herramienta esencial en la medicina con asistencia de inteligencia artificial \cite{ref35}.

Las Redes Neuronales Convolucionales (CNN) son arquitecturas de aprendizaje profundo elaboradas específicamente para procesar información dotada de estructura espacial, particularmente imágenes \cite{ref36}. Estas redes extraen de forma autónoma descriptores mediante operaciones de convolución, reducción de dimensionalidad (pooling) y transformaciones no lineales de activación, aprendiendo patrones que van desde representaciones simples hasta conceptos altamente abstractos \cite{ref36}. Debido a este mecanismo operacional, las CNN se emplean ampliamente en problemas de categorización, localización y delineación de imágenes \cite{ref36}. Su capacidad inherente para adquirir atributos visuales de manera automática las ha constituido como una de las arquitecturas más prominentes en campos como la visión computacional e imágenes médicas \cite{ref36}.

Los Vision Transformers (ViT) constituyen arquitecturas de aprendizaje profundo que transportan el mecanismo de auto-atención (self-attention) característico de Transformers al dominio del análisis de imágenes \cite{ref37}. Estos modelos segmentan una imagen en parches de tamaño uniforme que son procesados como secuencias para capturar tanto relaciones de proximidad como patrones globales \cite{ref37}. Derivado de este enfoque, los Vision Transformers se utilizan en categorización, localización y segmentación de imágenes \cite{ref37}. Su competencia en el modelado de dependencias de alcance extenso los ha posicionado como una alternativa equivalente a las redes neuronales convolucionales en aplicaciones de visión computacional \cite{ref37}.

Los modelos híbridos CNN-Transformer son arquitecturas que sintetizan las capacidades de redes neuronales convolucionales y Vision Transformers en el análisis de escenas visuales \cite{ref38}. Estos modelos aprovechan las CNN para la extracción de características de alcance reducido y los Transformers para la captura de relaciones de alcance global mediante mecanismos de auto-atención \cite{ref38}. Como resultado, producen representaciones más comprehensivas que mejoran substancialmente el desempeño en categorización, localización y delineación de imágenes \cite{ref38}. Gracias a este equilibrio entre exactitud y capacidad representativa, los modelos híbridos configuran una alternativa de considerable potencial para aplicaciones de visión computacional e imágenes médicas \cite{ref38}.

El aprendizaje por transferencia (Transfer Learning) constituye una metodología que reutiliza modelos previamente entrenados para resolver tareas novedosas empleando volúmenes de datos limitados \cite{ref39}. Este enfoque aprovecha el conocimiento acumulado por redes profundas para abreviar tiempos de entrenamiento y potenciar el rendimiento en nuevas aplicaciones \cite{ref39}. Adicionalmente, frecuentemente se combina con estrategias tales como expansión sintética de datos (data augmentation) y optimización paramétrica (fine-tuning) para incrementar la aptitud del modelo para generalizar a datos no observados \cite{ref39}. Gracias a estas ventajas operativas, el aprendizaje por transferencia goza de amplia utilización en sistemas de análisis y categorización de imágenes médicas basados en inteligencia artificial \cite{ref39}.

Las métricas de evaluación constituyen medidas numéricas empleadas para cuantificar el desempeño y la aptitud de generalización de un modelo de aprendizaje automático \cite{ref40}. La selección de una métrica particular depende de la naturaleza del problema de categorización y de los objetivos específicos de la aplicación \cite{ref40}. Entre las medidas más frecuentemente utilizadas figuran la exactitud (accuracy), la precisión (precision), la sensibilidad (recall), la puntuación F1 y el área bajo la curva ROC (AUC), siendo cada una de ellas indicativa de un aspecto distintivo del rendimiento del modelo \cite{ref40}. El análisis conjunto de estas métricas permite una valoración más íntegra y confiable del comportamiento de un sistema de aprendizaje automático \cite{ref40}.

La comparación de arquitecturas de aprendizaje profundo para categorización de imágenes médicas consiste en el análisis sistemático del desempeño de distintos modelos con el propósito de identificar cuál se adecua de forma óptima a una tarea determinada \cite{ref41}. Este proceso valora aspectos tales como la exactitud, la eficiencia computacional y la aptitud de generalización de arquitecturas fundadas en Redes Neuronales Convolucionales y Vision Transformers \cite{ref41}. La comparación sistemática permite identificar las ventajas inherentes y limitaciones de cada modelo en función de las particularidades del conjunto de datos y de la naturaleza del problema de categorización \cite{ref41}. Estos análisis constituyen una etapa cardinal en la selección de la arquitectura más apropiada en el desarrollo de sistemas de asistencia diagnóstica automatizada mediante inteligencia artificial \cite{ref41}.
